\subsection{COMPLEMENT OF A SET. THE EMPTY SET}

The relation \((x \notin A \text{ and } x \in X)\) is collectivizing in x by C51.

DEFINITION 3. Let A be a subset of a set X. The set of elements of X which do not belong to A, that is to say the set \(\mathfrak{E}\mathbf{x}\) ( \(x \notin \mathrm{A}\) and \(x \in \mathrm{X}\) ), is called the complement of A in X, and is denoted by \(\complement_{\mathbf{X}}\mathrm{A}\) or \(\mathrm{X} - \mathrm{A}\) (or by \(\complement\mathrm{A}\) if there is no risk of confusion).

Let A be a subset of a set X; the relations `` \(x \in X\) and \(x \notin A\) '' and \(x \in \complement_{X} A\) are then equivalent. Consequently the relation `` \(x \in X\) and \(x \notin \complement_{X} A\) '' is equivalent to `` \(x \in X\) and ( \(x \notin X\) or \(x \in A\) )'', hence to \(x \in A\) . In other words, \(A = \complement_{\mathbf{X}}(\complement_{\mathbf{X}} A)\) is a true relation. Similarly, one shows that if \(B\) is a subset of \(X\) , the relations \(A \subset B\) and \(\complement_{\mathbf{X}} B \subset \complement_{\mathbf{X}} A\) are equivalent.

THEOREM 1. The relation \((\forall x)(x \notin X)\) is functional in \(X\) .

For the relation \((\forall x)(x \notin X)\) implies \((\forall Y)(X \subset Y)\) ; by virtue of the axiom of extent, the relation \((\forall x)(x \notin X)\) is therefore single-valued in \(X\) . On the other hand, the relation \((\forall x)(x \notin C_{Y}Y)\) is true, which proves that \((\exists X)(\forall x)(x \notin X)\) is true.

The term \(\tau_x((\forall x)(x \notin X))\) corresponding to this functional relation is represented by the functional symbol \(\phi\) , and is called the empty set (*); the relation \((\forall x)(x \notin X)\) , which is equivalent to \(X = \phi\) , is read as follows: ``the set \(X\) is empty''. We have the theorems \(x \notin \phi\) , \(\phi \subset X\) , \(\complement_X X = \phi\) , \(\complement_X \phi = X\) . The relation \(X \subset \phi\) is equivalent to \(X = \phi\) . If \(R \{x\}\) is a relation, the relation \((\forall x)((x \in \phi) \Longrightarrow R \{x\})\) is true.

Remark. There exists no set of which every object is an element; in other words, ``not \((\exists X)(\forall x)(x \in X)\) '' is a theorem. For if there were such a set, then by C52 every relation would be collectivizing. But we have seen (no. 4) that the relation \(x \notin x\) is not collectivizing.

